\documentclass{article}
\usepackage{iclr2025_conference,times}

\usepackage{graphicx}
\usepackage{hyperref}
\usepackage{url}

\iclrfinalcopy
\pagestyle{empty}

\title{Adapting to Drift: Weather-Pattern Experts for Short-Term Photovoltaic Power Forecasting}
\author{\textbf{Wentao Wang$^{5}$, Haiyan Lu$^{1}$, Ayesha Ubaid$^{1}$, Fanyi Yang$^{2}$, Runyao Yu$^{3,4}$, Jianzhou Wang$^{5}$}\\[0.25em]\footnotesize $^1$University of Technology Sydney, $^2$Xidian University, $^3$Delft University of Technology\\\footnotesize $^4$Austrian Institute of Technology, $^5$Macau University of Science and Technology}

\makeatletter
\newcommand{\DeepPriorLogoRow}{%
  \AddToShipoutPictureFG*{%
    \AtPageUpperLeft{%
      \put(\LenToUnit{8.10in},\LenToUnit{-1.23in}){%
        \makebox[0pt][r]{
\includegraphics[width=0.825in,keepaspectratio]{deepprior-logo.png}}%
      }%
    }%
  }%
}

\newcommand{\DeepPriorTitleRow}{%
  \noindent\makebox[\textwidth][c]{%
    \begin{minipage}[c]{\textwidth}
      \centering\LARGE\bfseries \@title
    \end{minipage}%
  }\par
}

\newcommand{\DeepPriorAuthors}{%
  \begin{center}
    \begin{tabular}{c}\@author\end{tabular}
  \end{center}
}

\newcommand{\DeepPriorAbstract}[1]{%
  \begin{center}{\large\sc Abstract}\end{center}\vspace{-0.35em}%
  \begin{quote}#1\end{quote}%
}

\newcommand{\DeepPriorFigure}{%
  \noindent\makebox[\textwidth][c]{\includegraphics[width=\textwidth,keepaspectratio]{figure.png}}\par
  \vspace{2pt}%
  \noindent\makebox[\textwidth][c]{{\small Figure~1: Model Structure of ADrift.}}%
}

\makeatother
\begin{document}
\fancyhf{}
\renewcommand{\headrulewidth}{0pt}
\thispagestyle{empty}
\makeatletter
\DeepPriorLogoRow
\noindent\begin{minipage}[t][\dimexpr\textheight-\baselineskip\relax][s]{\textwidth}
\vspace*{0.50in}
\vfill
\DeepPriorTitleRow
\vfill
\DeepPriorAuthors
\vfill
\DeepPriorAbstract{Photovoltaic (PV) power forecasting is often developed under an implicit assumption of stationarity, yet the weather--power relationship evolves over time and induces distribution shifts commonly known as concept drift. Existing PV power forecasting models either ignore this issue entirely or address it in a limited and insufficient manner. To better understand this non-stationarity, we distinguish between internal and external drift, which motivate the key components of our design. Therefore, we present ADrift, a drift-aware forecasting framework that integrates patch-based temporal modeling, weather-guided representation learning, and lightweight online adaptation. To model internal drift, the backbone employs prototype-guided weather experts that capture diverse meteorological patterns within each input window. To cope with external drift, a learnable adapter updates its parameters through a temporal gap attention mechanism that enables targeted adjustments to the model. In addition, a proactive update strategy further mitigates supervision delays under rapidly changing conditions. Experiments on three real-world PV datasets show that ADrift consistently improves forecasting accuracy over static and online-learning baselines, demonstrating its potential for practical deployment under evolving weather conditions.}
\vfill
\DeepPriorFigure
\end{minipage}
\makeatother
\end{document}